\item[IELT（阴道内射精潜伏期）] 指阴茎插入阴道到射精的时间，是量化早泄的核心指标之一；规范测量采用秒表法，连续记录多次取中位数，以减少单次波动的影响。
\item[PEDT（早泄诊断工具）] 国际通用的 5 条目早泄自评量表，总分 0--20 分：≥11 分支持早泄诊断，9--10 分疑似，≤8 分不支持。
\item[自然变异性早泄] 射精过快呈偶发、情境性而非持续存在，控制力时好时坏，属正常范围内的波动，一般无需药物治疗。
\item[早泄样射精功能障碍] 主观自认"早泄"但射精潜伏期实际正常，困扰多源于对时长的过高期待或勃起困难时的"抢时间"心理，处理以心理教育与认知调整为主。
